% COMS 6998, Fall 2026 -- P1: Infrastructure & Baselines (ReasonLink)
% Same formatting as neural-symbolic-proposal.tex.
% Build: latexmk -pdf neural-symbolic-p1.tex
% Source of truth for every number here: PROGRESS.md in the repository.
\documentclass[conference,10pt,letterpaper]{IEEEtran}
\usepackage[T1]{fontenc}
\usepackage{cite}
\usepackage{amsmath,amssymb}
\usepackage{booktabs}
\usepackage{tabularx}
\usepackage{xcolor}
\usepackage{flushend}
\usepackage[hyphens]{url}
\usepackage[hidelinks]{hyperref}
\newcolumntype{Y}{>{\raggedright\arraybackslash}X}
\newcolumntype{P}[1]{>{\raggedright\arraybackslash}p{#1}}

% Red placeholder; swap the definitions to hide highlighting.
\newcommand{\fillme}[1]{\textcolor{red}{[#1]}}
% \newcommand{\fillme}[1]{#1}

\title{ReasonLink: P1 Infrastructure and Baselines}
\author{\IEEEauthorblockN{Vayun Gupta (vg2684), Aaron Mano Cherian (amc2535)}
\IEEEauthorblockA{COMS 6998: AI-Native Computing $\mid$ Fall 2026\\
P1 Milestone $\mid$ Due October 19, 2026, 11:59 PM ET}}

\begin{document}
\maketitle
\thispagestyle{plain}
\pagestyle{plain}

\noindent\textbf{Repository:} \url{https://github.com/AaronCherian404/ReasonLink}\\
All numbers below are reproducible from the repository (raw per-run data in \texttt{results/p1/}); \texttt{PROGRESS.md} logs every step, decision, and command.

% =====================================================================
\section{Summary and Platform Change}
% =====================================================================
We could not obtain the VCU118 that the Chipyard/FireSim flow needed, so we moved from Chipyard and Gemmini to ESP~\cite{esp}, the open SoC platform that already integrates NVDLA~\cite{nvdla} as a third-party accelerator alongside Ariane~\cite{espnvdla}. The neural accelerator is now NVDLA \emph{nv\_small} (int8, 64 MACs); the symbolic accelerator is still SAT-Accel~\cite{sataccel}. ESP's full flow is installed on a departmental server, and a VCU118 will be available later through ESP. The research question, the visual-Sudoku workload, the two baselines, and the four ReasonLink mechanisms are unchanged.

For P1 we built every piece of the pipeline that does not need the board, checked each against a reference (Table~\ref{tab:status}), and ran both baselines end to end (Sec.~\ref{sec:baselines}). The main gap is that SAT-Accel itself cannot run yet; a MiniSat stand-in consumes SAT-Accel's exact memory image instead (Sec.~\ref{sec:sat}).

\begin{table}[h]
\centering
\caption{Component status (October 16).}
\label{tab:status}
\footnotesize
\begin{tabularx}{\columnwidth}{@{}P{0.33\columnwidth}P{0.15\columnwidth}Y@{}}
\toprule
\textbf{Component} & \textbf{Status} & \textbf{Evidence} \\
\midrule
ESP toolchain on server & Ready & ESP 2026.1.0 pinned; Vivado, ModelSim, RISC-V GCC present \\
NVDLA compile flow & Working & ESP's loadable reproduced byte for byte \\
int8 LeNet classifier & Working (emulated) & 99.20\% MNIST; emulator matches NVDLA reference \\
Sudoku workload & Working & 1{,}000 puzzles + NVDLA output trace \\
SAT-Accel host image & Working & C port byte-exact vs.\ SAT-Accel, 13/13 \\
SAT solving & Stand-in & MiniSat on SAT-Accel's image; SAT-Accel blocked \\
Serial / double-buffered baselines & Working (host) & End to end, cycle-counted, 1{,}000 puzzles; RISC-V build \\
NVDLA, SAT-Accel on ESP SoC & Not started & After P1 \\
\bottomrule
\end{tabularx}
\end{table}

% =====================================================================
\section{Setup by Component}
% =====================================================================
\subsection{Repository}
All code is in one repository: \texttt{workload/} (CNF encoder, dataset, classifier, NVDLA compile script), \texttt{handoff/} (SAT-Accel image builder, solver backend, tests), \texttt{models/} (retrained weights), and \texttt{results/p1/} (machine-readable outputs). SAT-Accel, MiniSat, and Xilinx's open-source \texttt{ap\_int} headers are git submodules pinned to fixed commits. Local work runs in WSL Ubuntu 24.04 with a pinned Python environment (\texttt{env/setup\_wsl.sh}).

\subsection{ESP Platform (server)}
The server runs RHEL 8.10 (ESP's primary platform) with 16 cores and 62\,GB RAM, and already provides ESP's required tools: Vivado 2023.2, ModelSim DE 2023.2, Xcelium 23.09, Vitis HLS 2023.2, and RISC-V GCC 9.2, so ESP's 1--2\,h toolchain build is unnecessary. ESP is pinned to release 2026.1.0 (\texttt{a45f2bb}); only the NVDLA software submodule is fetched so far. The server lacks Vitis \texttt{v++} and XRT, which matters for SAT-Accel (Sec.~\ref{sec:sat}).

\subsection{NVDLA Classifier}
We use the LeNet~\cite{lenet} that ESP ships for its NVDLA demo, so our classifier is the network ESP already runs on NVDLA.
\begin{itemize}
\item \emph{Compiler.} Recompiling ESP's model with the prebuilt \texttt{nvdla\_compiler} reproduces ESP's shipped loadable byte for byte with \texttt{-{}-quantizationMode per-filter -{}-informat nchw}; the compiler is deterministic, and no other flag combination matched.
\item \emph{int8 emulator.} A Python model of the int8 pipeline (the runtime's input mapping, per-filter weight scales, calibrated requantization) reproduces ESP's published NVDLA output for its test image exactly. It generates the per-cell scores that stand in for NVDLA's output buffer.
\item \emph{Weights.} ESP's demo weights reach only 74.5\% on MNIST. They were trained with mean-image subtraction and a 1/80 input scale that neither the deploy model nor the NVDLA int8 input path applies. We retrained the same architecture on the input NVDLA actually receives and wrote a new calibration table (Table~\ref{tab:nn}). The retrained model compiles with ESP's flags to a 445{,}736-byte loadable.
\end{itemize}

\begin{table}[h]
\centering
\caption{LeNet accuracy on the MNIST test set (10{,}000 images).}
\label{tab:nn}
\footnotesize
\begin{tabularx}{\columnwidth}{@{}Yrr@{}}
\toprule
\textbf{Weights} & \textbf{float} & \textbf{int8 (emulated)} \\
\midrule
ESP demo, under NVDLA preprocessing & 74.5\% & 74.5\% \\
Retrained (ours) & 99.17\% & 99.20\% \\
\bottomrule
\end{tabularx}
\end{table}

\subsection{Workload}
Sudoku uses the standard encoding with variable $81r + 9c + d$: 729 variables and 11{,}988 base clauses, identical for every puzzle; predicted givens are added as unit clauses. We generated 1{,}000 puzzles with unique solutions and 24--40 clues (mean 32.1), in the style of visual Sudoku~\cite{satnet}. Each given cell holds an MNIST \emph{test} image of its digit. The classifier's int8 logits for all puzzles are stored as a binary trace, the input to every interface.

\subsection{SAT Solver}
\label{sec:sat}
\emph{Host conversion.} SAT-Accel's host code builds a device memory image from the CNF: paged clause and occurrence stores, 64-byte per-variable metadata, and a decision stack. We ported this conversion to standalone C, since it is exactly the work the serial baseline must time. A golden test compiles SAT-Accel's own unmodified parser next to our port and compares every region byte for byte; all 13 instances match: four Sudoku instances (one of them UNSAT) and nine of SAT-Accel's benchmarks, up to 758 variables and 47{,}820 clauses.

\emph{Solving.} SAT-Accel's simulation flow requires Vitis \texttt{v++} and XRT, which the server lacks and which need root to install. For P1 a MiniSat~\cite{minisat} backend reads the clauses back out of the SAT-Accel image, as the device's copy-in would, and solves them. A PySAT~\cite{pysat} run serves as the software reference. Later, the real accelerator replaces the backend behind the same interface.

% =====================================================================
\section{Results So Far}
% =====================================================================
\textbf{Functional pipeline.} On all 1{,}000 puzzles, the C path (trace $\rightarrow$ argmax $\rightarrow$ SAT-Accel image $\rightarrow$ MiniSat) agrees with the PySAT reference on SAT/UNSAT for every puzzle. It returns the identical grid for 992; the other eight have several solutions because of one misread given, and we verified that both grids satisfy every clause. Table~\ref{tab:func} summarizes accuracy.

\begin{table}[h]
\centering
\caption{Visual Sudoku, 1{,}000 puzzles.}
\label{tab:func}
\footnotesize
\begin{tabularx}{\columnwidth}{@{}Yr@{}}
\toprule
\textbf{Metric} & \textbf{Value} \\
\midrule
Given-cell accuracy (int8) & 99.26\% \\
Puzzles solved correctly & 79.0\% \\
UNSAT from misread givens & 19.2\% \\
SAT but wrong & 1.8\% \\
Solver work per puzzle (decisions / conflicts) & 1.2 / 0.16 \\
\bottomrule
\end{tabularx}
\end{table}

\textbf{Handoff size.} Table~\ref{tab:bytes} gives bytes per puzzle on the native SAT-Accel path. The neural stage produces 1.3\,KB of logits that encode 36 givens. SAT-Accel's host flow nonetheless rebuilds a 632\,KB image, almost all of it the unchanging base formula, and migrates its full 10.5\,MB of allocated buffers on every instance.

\textbf{Baselines.}
\label{sec:baselines}
Both baselines run end to end on all 1{,}000 puzzles. A host thread performs SAT-Accel's native flow: argmax, image build, copy into separate device buffers, and a doorbell plus completion poll. A device thread decodes its copy and solves it. Every phase is timed with the cycle counter (\texttt{rdtsc}; \texttt{rdcycle} on RISC-V). The double-buffered baseline uses two buffer sets with explicit ownership states. We run each baseline \emph{faithful} (full buffers cleared and copied, as SAT-Accel's host does) and \emph{lean} (only the bytes in use), so that the serial baseline is not a strawman. All four runs return identical answers that match the PySAT reference. Table~\ref{tab:baselines} gives the breakdown.

\begin{table}[h]
\centering
\caption{Baseline breakdown, median $\mu$s per puzzle (host x86, 1{,}000 puzzles). Share: $T_{\text{handoff}}/(T_{\text{handoff}}+T_{\text{sat}})$ for serial runs; exposed fraction of wall time for double-buffered runs.}
\label{tab:baselines}
\footnotesize
\setlength{\tabcolsep}{3pt}
\begin{tabularx}{\columnwidth}{@{}Yrrrrr@{}}
\toprule
\textbf{Run} & $T_{\text{conv}}$ & $T_{\text{xfer}}$ & $T_{\text{sat}}$ & \textbf{Share} & \textbf{Wall} \\
\midrule
Serial, faithful & 1112 & 880 & 728 & 73.6\% & 2809 \\
Serial, lean & 696 & 13 & 726 & 50.5\% & 1450 \\
Double, faithful & 1461 & 1159 & 764 & 70.5\% & 2703 \\
Double, lean & 711 & 16 & 743 & 5.1\% & 801 \\
\bottomrule
\end{tabularx}
\end{table}

These pilot numbers support H1 and H2. In the native flow, the handoff is 74\% of end-to-end time, yet the argmax that actually consumes the neural output takes under 1\,$\mu$s; the rest is rebuilding and moving the base formula. Double buffering hides the handoff only when it is shorter than the solve. The faithful handoff (2.6\,ms vs.\ a 0.76\,ms solve) stays 70\% exposed; the lean one is nearly hidden but still moves 632\,KB per puzzle. $T_{\text{sync}}$ is below 1\,$\mu$s with busy polling, whereas SAT-Accel's host polls once per second.

\begin{table}[h]
\centering
\caption{Data per puzzle on the native SAT-Accel path (puzzle 0, 36 givens).}
\label{tab:bytes}
\footnotesize
\begin{tabularx}{\columnwidth}{@{}Yr@{}}
\toprule
\textbf{Quantity} & \textbf{Bytes} \\
\midrule
NVDLA output (int8 logits, $81 \times 16$\,B) & 1{,}296 \\
SAT-Accel image, bytes in use & 632{,}436 \\
SAT-Accel buffers migrated per instance & 10{,}536{,}436 \\
\bottomrule
\end{tabularx}
\end{table}

\textbf{Limitations.} All times are host x86 cycles; on ESP, conversion runs on Ariane and transfers use DMA over the NoC, so absolute values and ratios will change. The same harness cross-compiles to a static RISC-V binary for that purpose. Solver time in P1 is MiniSat CPU time, not SAT-Accel cycles. NVDLA has not yet executed on hardware: the loadable is compiled, and its outputs are emulated. NVDLA's real output layout in memory is still to be confirmed on ESP. Sudoku with $\geq$24 clues is nearly propagation-only, so the H3 crossover will need harder instances.

% =====================================================================
\section{Next Steps}
% =====================================================================
\label{sec:next}
\begin{center}
\footnotesize
\begin{tabularx}{\columnwidth}{@{}P{0.15\columnwidth}Y@{}}
\toprule
\textbf{Target} & \textbf{Task} \\
\midrule
Nov 2 (P2) & Run SAT-Accel functionally (its HLS kernels as threads, or XRT on a machine with access); first ESP SoC with Ariane and NVDLA in simulation, running the RISC-V baselines; ReasonLink prototype with literal stream and on-chip FIFO, compared against Table~\ref{tab:baselines}. \\
After P2 & Delta transfers and overlap; SAT-Accel on ESP; VCU118 runs when the board is available. \\
\bottomrule
\end{tabularx}
\end{center}

\noindent\textbf{Team contributions.} \fillme{Who did what for P1.}

\bibliographystyle{IEEEtran}
\bibliography{neural-symbolic-references,p1-additional-references}
\end{document}
